% --------------------------------------------------------- INTRODUCTION
\section{Introduction}

Crystal discovery is shifting from generating candidate structures to deciding which
predictions warrant calculation and experiment. High-throughput databases, crystal
generators and tool-using agents
\cite{jain2013commentary,merchant2023scaling,xie2022crystal,jiao2023crystal,sriram2024flowllm,antunes2024crystal,boiko2023autonomous,mbran2024augmenting}
now supply candidates faster than experiment or density functional theory (DFT)
calculations of thermodynamic energies and phonon spectra can assess them
\cite{bartel2019new,petretto2018highthroughput,zhu2024highthroughput,szymanski2023autonomous}.
Yet many generative pipelines test little more than a fixed
minimum interatomic distance
\cite{court20203d,xie2022crystal,jiao2023crystal,miller2024flowmm,betala2025lematgenbench},
and avoiding gross overlap does not make coordination, electrostatics,
bond valence or chemical ordering plausible. A recent Comment argues that data alone will struggle to deliver materials discovery, and that AI
must learn the chemical rules governing atomic arrangements \cite{smit2026dataonly}. What
is missing is a rapid, interpretable set of laws for structural plausibility that identify
the physical or chemical constraint a structure violates. Such laws would sit between elementary geometry and the
costlier questions of thermodynamic stability, dynamical stability and synthesizability.

Pauling's five rules offer a richer precedent. They connect ionic size, coordination
and electrostatic valence, and they favour simpler structures. Requiring only a radius
table and formal charges, they exemplify a tradition of simple laws whose chemical
reasoning can be examined directly
\cite{pauling1929principles,goldschmidt1926gesetze,pettifor1984chemical,zunger1980systematization,hawthorne2014structure}.
But these rules and the modern distance cutoff err in opposite directions. An audit found
that only 13\% of about 5{,}000 oxides satisfied rules 2--5 together
\cite{george2020limited}, and in our evaluations only 6.5\% of charge-balanced ionic
experimental structures did. At the opposite extreme, the 0.5- and 0.7-\AA{}
minimum-distance cutoffs used in generative pipelines detected just 1.6\% and 3.2\% of
chemically damaged structures. Applied jointly, the classical
rules are too idealised, and the distance cutoff is inexpensive but element-blind.
Experimental-structure satisfaction alone cannot show what a criterion rejects. A good law
must therefore keep experimental
structures, detect damaged ones at a high rate and say why they are implausible.

Two tests determine the value of such laws: whether they track synthesizability and whether
they identify structural errors. Databases record many successes but few definitive
failures, so synthesizability has been estimated using crystal-likeness scores
\cite{jang2020structure}, language models \cite{song2025llm} or recommendation engines
\cite{griesemer2025wideranging}. None of them tests an explicit chemical law, so it remains unknown whether a structural
criterion fitted without synthesis labels tracks synthesizability.
The second test asks whether the structure handed to DFT or experiment is correct
at all. Every later assessment assumes it is, yet prominent cases show that the chemistry
can be misassigned.
GNoME reported 381{,}000 newly discovered structures on the convex hull, bringing its stable
set to 421{,}000 \cite{merchant2023scaling}. This set remains enriched in rare
low-symmetry structures that pass distance cutoffs. A perspective traced part of the excess
to the artificial ordering of similar elements, such as rare-earth pairs and Zr--Hf,
over sites that should be equivalent \cite{cheetham2024artificial,smit2026dataonly}. Chemical
ordering also shaped A-Lab's success criterion, which accepted ordered and partially
disordered products \cite{szymanski2023autonomous}. A reanalysis disputed several novelty
and phase assignments \cite{leeman2024challenges}, and an author correction later found
four of A-Lab's 40 successes inconclusive \cite{szymanski2026correction}.
A more severe failure places the wrong element on a single site. Harrison and colleagues
documented at least 70 falsified structures in which genuine diffraction data were paired
with altered element identities \cite{harrison2010falsified}. An archive of retracted
depositions records the same pattern across Cu, Ni, Mn and Fe labels
\cite{iucr2012retraction}. Distance cutoffs and energy calculations examine neither
chemical ordering nor elemental identity. That blind spot leaves two open cases: whether generated catalogues artificially order
similar elements, and how a wrong occupant at plausible coordinates can pass existing
structural screens. Closing either requires laws that name the physical constraint that fails.

Here we show that autonomous agents can discover such laws. Our agents received experimental
structures \cite{zagorac2019recent,grazulis2009crystallography}, interpretable quantities
and prescribed ways to damage a crystal. They generated, implemented and actively tested
candidate laws. Because failed claims remained on record, refutation, not proposal
volume, measured progress. Of two million candidate laws, eight survived: the Plausibility Rules
for Inorganic Structures (PRIS). Each law encodes one of five mechanisms: short-range
repulsion, ionic contact and packing, electrostatic balance, bond-valence conservation and
crystallographic site complexity. Every violation therefore identifies the mechanism that
fails. Sets of these laws reached 82--99\% experimental-structure satisfaction, and the
strictest set detected 87.9\% of chemically damaged structures.
No synthesis label entered the selection of the PRIS laws, yet PRIS plausibility is
linearly correlated with synthesizability. Our agents' PRIS-derived synthesis score (PSS)
screened 83.7\% of hard-to-synthesize structures while keeping 80.7\% of experimental ones,
and every removal was traceable to a named mechanism. In an inverse-design task targeting a
high bulk modulus, explainable screening by PRIS and PSS reduced the DFT validation queue for
MatterGen outputs by up to 67.3\%. That screen kept 99.2\% of the
structures that reached the design target under DFT. The same mechanisms close both open cases: they trace GNoME's low-symmetry excess to
artificial ordering and expose wrong-element assignments hidden behind plausible coordinates. PRIS thereby turns crystal screening from a
pass-or-fail check into a chemical diagnosis. Autonomous agents required to refute their own claims can discover physicochemical laws
that guide the calculations and experiments that follow.
